
\paragraph{\SGARefLink{sga3:VIIB:5.2.1.A}{5.2.1.A}.}
\phantomsection
\label{sga3:VIIB:5.2.1.A}
\footnote{\label{sga3:VIIB:nde:155}The editors added
\SGARefLink{sga3:VIIB:5.2.1.A}{5.2.1.A} and \SGARefLink{sga3:VIIB:5.2.1.B}{5.2.1.B}.}
For every \(r\in\mathbb N\), put
\[
  \mathcal I_r=\overline{I^2}+I_r,
\]
so that
\[
  \mathcal I_r/\overline{I^2}\simeq\omega_r.
\]
There is then a commutative diagram with exact rows:
\begin{center}
  \phantomsection
  \label{sga3:VIIB:5.2.1.A:exact-rows}
  \begin{tikzcd}[row sep=3em,column sep=2.6em]
0 \arrow[r] & \mathcal I_{r-1} \arrow[r] \arrow[d]
  & \mathcal I_r \arrow[r] \arrow[d]
  & \mathcal I_r/\mathcal I_{r-1}
      \arrow[r] \arrow[d,"\wr"]
  & 0 \\
0 \arrow[r] & \omega_{r-1} \arrow[r]
  & \omega_r \arrow[r]
  & \omega_r/\omega_{r-1} \arrow[r]
  & 0
\end{tikzcd}

\end{center}

Since \(k\) is a field, the rows split. Thus a pseudobasis
\(\mathscr B_{r-1}\) of \(\mathcal I_{r-1}\) can be completed
to a pseudobasis
\[
  \mathscr B_{r-1}\cup\mathscr B'_r
\]
of \(\mathcal I_r\). The closed subspace \(S_r\) with
pseudobasis \(\mathscr B'_r\) is then a complement of
\(\mathcal I_{r-1}\) in \(\mathcal I_r\), and the projection
\[
  \pi:I\longrightarrow\omega
\]
induces an isomorphism from \(S_r\) onto a complement
\(\omega'_r\) of \(\omega_{r-1}\) in \(\omega_r\).

Let \(\mathcal I_\infty\) be the closed ideal
\[
  \mathcal I_\infty
  =
  \overline{\bigcup_r\mathcal I_r}.
\]
It likewise has a complement \(S_\infty\) in \(I\), and
\(\pi\) induces an isomorphism from
\[
  \mathcal I_\infty/\overline{I^2}
\]
onto the closure
\[
  \omega_\infty
  =
  \overline{\bigcup_r\omega_r}
\]
of the union of the \(\omega_r\), and an isomorphism from
\(S_\infty\) onto a complement \(\omega'_\infty\) of
\(\omega_\infty\) in \(\omega\). Denote the latter isomorphism by
\[
  \eta:S_\infty\xrightarrow{\sim}\omega'_\infty.
\]
We thus obtain continuous linear maps:
\begin{center}
  \phantomsection
  \label{sga3:VIIB:5.2.1.A:splitting-maps}
  \begin{tikzcd}[row sep=3.2em,column sep=3em]
\overline{I^2}\times S_\infty\times\widehat{\bigoplus}_r S_r
  \arrow[r,"\phi"]
  \arrow[d,"\eta" left]
  \arrow[d,phantom,"{\times\theta}" right]
  & I \\
\omega=\omega'_\infty\times\omega_\infty & {}
\end{tikzcd}

\end{center}

Here
\[
  \widehat{\bigoplus}_r S_r
\]
is the direct sum of the \(S_r\) in \(\operatorname{PC}(k)\),
that is,
\[
  \left(\prod_r S_r^\dagger\right)^*
\]
\textup{(cf.\ editorial note \textup{(16)} to \SGARefLink{sga3:VIIB:0.2.2}{0.2.2})}, and
\[
  \theta:
  \widehat{\bigoplus}_r S_r
  \longrightarrow\omega_\infty
\]
is induced by the maps
\[
  S_r\xrightarrow{\sim}\omega'_r
  \lhook\joinrel\longrightarrow\omega_\infty.
\]
It follows that a sufficient condition \textup{(but not a
necessary one, as seen below)} for obtaining a section
\[
  \sigma:\omega\longrightarrow I
\]
of the desired kind is that \(\theta\) be an isomorphism.
By duality \textup{(cf.\ \SGARefLink{sga3:VIIB:0.2.2}{0.2.2})}, this is equivalent to saying
that the linear map
\[
  \omega_\infty^\dagger
  \longrightarrow
  \prod_r S_r^\dagger
\]
is bijective.
